%% start of file `template.tex'.
%% Copyright 2006-2015 Xavier Danaux (xdanaux@gmail.com), 2020-2022 moderncv maintainers (github.com/moderncv).
%
% This work may be distributed and/or modified under the
% conditions of the LaTeX Project Public License version 1.3c,
% available at http://www.latex-project.org/lppl/.


\documentclass[11pt,letterpaper,sans]{moderncv}        % possible options include font size ('10pt', '11pt' and '12pt'), paper size ('a4paper', 'letterpaper', 'a5paper', 'legalpaper', 'executivepaper' and 'landscape') and font family ('sans' and 'roman')

\moderncvcolor{black}                               % color options 'black', 'blue' (default), 'burgundy', 'green', 'grey', 'orange', 'purple' and 'red'
% moderncv themes
\moderncvstyle{classic}                             % style options are 'casual' (default), 'classic', 'banking', 'oldstyle' and 'fancy'
%\renewcommand{\familydefault}{\sfdefault}         % to set the default font; use '\sfdefault' for the default sans serif font, '\rmdefault' for the default roman one, or any tex font name
%\nopagenumbers{}                                  % uncomment to suppress automatic page numbering for CVs longer than one page

% adjust the page margins
\usepackage[scale=0.75]{geometry}
%\setlength{\footskip}{136.00005pt}                 % depending on the amount of information in the footer, you need to change this value. comment this line out and set it to the size given in the warning
\setlength{\hintscolumnwidth}{3cm}                % if you want to change the width of the column with the dates
%\setlength{\makecvheadnamewidth}{10cm}            % for the 'classic' style, if you want to force the width allocated to your name and avoid line breaks. be careful though, the length is normally calculated to avoid any overlap with your personal info; use this at your own typographical risks...

% font loading
% for luatex and xetex, do not use inputenc and fontenc
% see https://tex.stackexchange.com/a/496643
\ifxetexorluatex
  \usepackage{fontspec}
  \usepackage{unicode-math}
  \defaultfontfeatures{Ligatures=TeX}
  \setmainfont{Latin Modern Roman}
  \setsansfont{Latin Modern Sans}
  \setmonofont{Latin Modern Mono}
  \setmathfont{Latin Modern Math} 
\else
  \usepackage[T1]{fontenc}
  \usepackage{lmodern}
\fi

% document language
\usepackage[french]{babel}  % FIXME: using spanish breaks moderncv

% personal data
\name{Annabelle Dussault}{}
% nom en noir
\renewcommand*{\firstnamestyle}[1]{{\namefont\textcolor{black}{#1}}}
\renewcommand*{\lastnamestyle}[1]{{\namefont\textcolor{black}{#1}}}
\title{Étudiante à la maîtrise en science politique}                               % optional, remove / comment the line if not wanted
%\born{4 July 1776}                                 % optional, remove / comment the line if not wanted
%\address{street and number}{postcode city}{country}% optional, remove / comment the line if not wanted; the "postcode city" and "country" arguments can be omitted or provided empty
%\phone[mobile]{+1~(234)~567~890}                   % optional, remove / comment the line if not wanted; the optional "type" of the phone can be "mobile" (default), "fixed" or "fax"
%\phone[fixed]{+2~(345)~678~901}
%\phone[fax]{+3~(456)~789~012}
\email{anna5.dussault@gmail.com}                               % optional, remove / comment the line if not wanted
\homepage{annabelle-dussault.github.io}                        % optional, remove / comment the line if not wanted

% Social icons
\social[linkedin]{annabelle-dussault-625a05303}                        % optional, remove / comment the line if not wanted
%\social[xing]{john\_doe}                           % optional, remove / comment the line if not wanted
%\social[twitter]{@anna5dussault}                             % optional, remove / comment the line if not wanted
%\social[github]{annabelle-dussault}                              % optional, remove / comment the line if not wanted
%\social[gitlab]{jdoe}                              % optional, remove / comment the line if not wanted
% \social[stackoverflow]{0000000/johndoe}            % optional, remove / comment the line if not wanted
% \social[bitbucket]{jdoe}                           % optional, remove / comment the line if not wanted
% \social[skype]{jdoe}                               % optional, remove / comment the line if not wanted
% \social[orcid]{0000-0000-000-000}                  % optional, remove / comment the line if not wanted
% \social[researchgate]{jdoe}                        % optional, remove / comment the line if not wanted
% \social[researcherid]{jdoe}                        % optional, remove / comment the line if not wanted
% \social[telegram]{jdoe}                            % optional, remove / comment the line if not wanted
% \social[whatsapp]{12345678901}                     % optional, remove / comment the line if not wanted
% \social[signal]{12345678901}                       % optional, remove / comment the line if not wanted
% \social[matrix]{@johndoe:matrix.org}               % optional, remove / comment the line if not wanted
% \social[googlescholar]{googlescholarid}            % optional, remove / comment the line if not wanted


%\extrainfo{additional information}                 % optional, remove / comment the line if not wanted
%\photo[64pt][0.4pt]{picture}                       % optional, remove / comment the line if not wanted; '64pt' is the height the picture must be resized to, 0.4pt is the thickness of the frame around it (put it to 0pt for no frame) and 'picture' is the name of the picture file
%\quote{Some quote}                                 % optional, remove / comment the line if not wanted

% bibliography adjustments (only useful if you make citations in your resume, or print a list of publications using BibTeX)
%   to show numerical labels in the bibliography (default is to show no labels)
%\makeatletter\renewcommand*{\bibliographyitemlabel}{\@biblabel{\arabic{enumiv}}}\makeatother
\renewcommand*{\bibliographyitemlabel}{[\arabic{enumiv}]}
%   to redefine the bibliography heading string ("Publications")
%\renewcommand{\refname}{Articles}

% bibliography with mutiple entries
%\usepackage{multibib}
%\newcites{book,misc}{{Books},{Others}}
%----------------------------------------------------------------------------------
%            content
%----------------------------------------------------------------------------------
\begin{document}
%\begin{CJK*}{UTF8}{gbsn}                          % to typeset your resume in Chinese using CJK
%-----       resume       ---------------------------------------------------------
\makecvtitle

\section{Formation}
\cventry{2026--auj.}{Maîtrise en science politique (M.A.)}{Université Laval}{Québec}{}{Avec mémoire. Direction : Yannick Dufresne}  % arguments 3 to 6 can be left empty
\cventry{2023--2026}{Baccalauréat en droit (LL.B.)}{Université Laval}{Québec}{}{}
\cventry{2021--2023}{Diplôme d'études collégiales en Sciences, lettres et arts}{Cégep de Sainte-Foy}{Québec}{}{}

%\section{Master thesis}
%\cvitem{title}{\emph{Title}}
%\cvitem{supervisors}{Supervisors}
%\cvitem{description}{Short thesis abstract}

\section{Expérience en recherche}
\cventry{2026--auj.}{Auxiliaire de recherche}{Chaire de leadership en enseignement des sciences sociales numériques (CLESSN), Université Laval}{Québec}{}{}
\cventry{2026--auj.}{Auxiliaire de recherche}{Chaire du Canada en économie politique internationale, Université Laval}{Québec}{}{}
\cventry{2024--2025}{Auxiliaire de recherche}{Chaire de leadership en enseignement René-Dussault sur l’inclusion des traditions autochtones dans les programmes de formation en droit, Université Laval}{Québec}{}{}

\section{Prix et distinctions}
\cvitem{2024}{Bourse d'excellence à l'admission (1er cycle), Faculté de droit, Université Laval}
\cvitem{2023}{Médaille du lieutenant-gouverneur pour la jeunesse, Lieutenant-gouverneur du Québec}
\cvitem{2020}{Premier prix, concours d'écriture historique, Société historique de Québec}
\cvitem{2020}{Bourse pour la qualité du français, Association pour le soutien et l'usage de la langue française (ASULF)}

\section{Groupes de recherche}
\cvitem{Depuis 2026}{Membre du Centre d'analyse des politiques publiques (CAPP), Université Laval}
\cvitem{Depuis 2026}{Membre du Centre interdisciplinaire de recherche sur la gouvernance mondiale (CIRGoM), Université Laval}
\cvitem{Depuis 2026}{Membre du Centre pour l'étude de la citoyenneté démocratique (CECD), Université Laval}

\newpage
\section{Expérience professionnelle}
\cventry{2025--2026}{Page}{Assemblée nationale du Québec}{Québec}{}{}
\cventry{2025}{Recherchiste-étudiante}{Cour d'appel du Québec}{Québec}{}{}

\section{Engagement et activités parascolaires}
\cventry{2023--auj.}{Parlement jeunesse du Québec (PJQ)}{}{}{}{
\begin{itemize}
\item Ministre des Institutions démocratiques (2026)
\item Présidente de la Commission des institutions (2025)
\item Déléguée internationale auprès du Conseil des jeunes valdôtains, Italie (2025)
\item Responsable de dossier au journal La Colline (2024)
\item Députée et vice-présidente de la Commission des relations avec les citoyens (2023)
\end{itemize}}
\cventry{2026--2027}{Administratrice responsable des partenariats étrangers}{Association québécoise des jeunes parlementaires (AQJP)}{}{}{}
\cventry{2026--2027}{Vice-présidente du comité des mini-simulations}{Association québécoise des jeunes parlementaires (AQJP)}{}{}{}
\cventry{2026}{Bénévole-animatrice}{Forum Québec PSP 2026 (Paix, Sécurité, Prospérité)}{}{}{}
\cventry{2024--2025}{Bénévole en recherche juridique}{Étudiants Pro Bono Canada, Centre de pédiatrie sociale de Lévis}{}{}{}
\cventry{2024--2025}{Vice-présidente aux affaires socioculturelles}{Simulation francophone des Nations Unies de l'Université Laval (SiFNUUL)}{}{}{}
\cventry{2022--2023}{Représentante des étudiants}{Groupe de travail et de réflexion sur la liberté académique, Cégep de Sainte-Foy}{}{}{}
\cventry{2022--2023}{Membre du jury local}{Prix littéraire des collégiens}{}{}{}

\section{Compétences}
\cvitem{Outils}{R (tidyverse, ggplot2), \LaTeX, Git et GitHub, Zotero}
\cvitem{Langues}{Français (langue maternelle), anglais (avancé)}
\end{document}


%% end of file `template.tex'.

